\begin{abstract}
Video diffusion transformers produce high-quality videos at substantial computational cost, limiting interactive and large-scale serving.
Existing acceleration methods typically optimize each request independently through feature caching, token pruning, or sparse attention, leaving redundancy across related generation requests largely unexplored.
We present \emph{Carnator}, a cross-request acceleration framework that reuses historical computation according to generation-level compatibility.
\emph{Carnator} first performs a lightweight early probe and constructs a generation-native object signature from early diffusion states, enabling compatibility-aware cache retrieval beyond prompt similarity.
A risk-constrained calibration procedure determines compatibility thresholds according to the observed quality risk of actual cache reuse.
Accepted caches are then exploited through joint latent-attention inference, which combines region-restricted target computation with historical sparse-connectivity reuse across complementary Transformer layer groups.
Experiments on diverse text-to-video requests show that CARNATOR improves reuse prediction by [X], reduces erroneous reuse by [Y], and achieves [Z]$\times$ end-to-end acceleration with limited quality degradation.
\end{abstract}

% \begin{abstract}
% Video diffusion transformers produce high-quality videos at substantial computational cost, limiting interactive and large-scale serving.
% Existing acceleration methods typically optimize each request independently through fewer denoising steps, feature caching, token pruning, or sparse attention, leaving cross-request redundancy largely unexplored.
% We present CARNATOR, a framework that accelerates video diffusion by reusing computation across related requests.
% CARNATOR performs a lightweight early probe and constructs a latent-attention signature from latent dynamics, temporal statistics, cross-attention entity responses, self-attention affinities, and attention-pattern statistics.
% The signature estimates compatibility between a new request and cached states at different anchor steps, identifies reusable regions, and selects an appropriate reuse mode.
% CARNATOR then performs hierarchical cache reuse: it reuses cached latents and intermediate features outside active regions while recomputing target-specific tokens, and reuses validated sparse-attention patterns within active regions to skip low-value query-key interactions.
% Safety connections and dense fallback preserve spatial and temporal consistency.
% Experiments on diverse text-to-video requests show that CARNATOR improves reuse prediction by [X], reduces erroneous reuse by [Y], and achieves [Z]x end-to-end acceleration with limited quality degradation.
% \end{abstract}